%------------------------------------------------------------------------------%
\begin{question}
  Dado \(z=\frac{ \left(1-\sqrt{3}\right) + \left(1+\sqrt{3}\right)i}{2+2i} \)

  calcule

  a) a parte real de $z$

  b) a parte imaginária de $z$

  c) o módulo de $z$

  d) o argumento de $z$
\end{question}

%------------------------------------------------------------------------------%
\begin{resolution}{Avaliando $z$}
\begin{align*}
  z
  & = \frac{ \left(1-\sqrt{3}\right) + \left(1+\sqrt{3}\right)i}{2+2i} \\
  & = \frac{ \left[\left(1-\sqrt{3}\right) + \left(1+\sqrt{3}\right)i\right](2-2i)}{(2+2i)(2-2i)} \\
  & = \frac{ \left(1-\sqrt{3}\right)(2-2i) + \left(1+\sqrt{3}\right)(2i+2)}{2^2 + 2^2} \\
  & = \frac{ 2-2i-2\sqrt{3}+2\sqrt{3}i + 2i+2+2\sqrt{3}i +2\sqrt{3}}{8} \\
  & = \frac{ 4 + 4\sqrt{3}i}{8}
  \;=\; \frac{1}{2} + \frac{\sqrt{3}}{2}i
\end{align*}
\end{question}

%------------------------------------------------------------------------------%
\begin{resolution}{Partes real e imaginária de $z$}
  Parte real de $z$
  \[
    \Re(z) = \frac{1}{2}
  \]

  Parte imaginária de $z$
  \[
    \Im(z) = \frac{\sqrt{3}}{2}
  \]
\end{resolution}

%------------------------------------------------------------------------------%
\begin{resolution}{Módulo de $z$}
  \[
    \abs{z}
    = \sqrt{\left(\frac{1}{2}\right)^2 + \left(\frac{\sqrt{3}}{2}\right)^2}
    = \sqrt{\frac{1}{4} + \frac{3}{4}}
    = \sqrt{\frac{1+3}{4}}
    = \sqrt{1}
    = 1
  \]
\end{resolution}

%------------------------------------------------------------------------------%
\begin{resolution}{Argumento de $z$}
  \[
    \varphi=\arg(z)
  \]
  \[
    \sin(\varphi)
    = \frac{\Im(z)}{\abs{z}}
    = \frac{\sqrt{3}}{2}
  \]
  \[
    \cos(\varphi)
    = \frac{\Re(z)}{\abs{z}}
    = \frac{1}{2}
  \]
  \[
    \arg(z)
    = \ang{60}
    = \frac{\pi}{3} \text{rad}
  \]
\end{resolution}

%------------------------------------------------------------------------------%




